\documentclass[10pt,a4paper]{amsart}
\usepackage[margin=19mm]{geometry}
\usepackage{amsmath,amssymb,mathtools}
\usepackage[scr=rsfs,cal=euler]{mathalfa}
\usepackage{xfrac}
\usepackage[unicode,hidelinks,bookmarks=false]{hyperref}
\newcommand{\ie}{i.e.\ }
\newcommand{\cf}{cf.\ }
\newcommand{\resp}{respectively\ }
\newcommand{\Subsection}{\subsection}
\providecommand{\nobreakdash}{\nobreak\hbox{-}}
\setlength{\emergencystretch}{2em}
\multlinegap0pt
\usepackage{mathtools}
\usepackage{amssymb}
\newtheorem{prop}{Proposition}
 \def\cI{\mathcal{I}}
\def\cK{\mathcal{K}}
 \def\cL{\mathcal{L}}
\newcommand{\commut}[2]{[#1{,}\,#2]}
\newcommand{\ffrac}[2]{\raisebox{.5pt}%
  {\footnotesize$\displaystyle\frac{#1}{#2}$}\kern1pt}
\newcommand{\half}{\mathchoice{%
    \ffrac{1}{2}}{\frac{1}{2}}{\frac{1}{2}}{\frac{1}{2}}}
\renewcommand{\tilde}{\widetilde}
\title{Background fields in the presymplectic BV-AKSZ approach}
\author{Ivan Dneprov, Maxim Grigoriev}
\date{Source: arXiv:2505.09885v2 (CC BY 4.0)}
\begin{document}
\maketitle

\noindent\emph{Source and licence.} Background fields in the presymplectic BV-AKSZ approach. Version arXiv:2505.09885v2 (CC BY 4.0), distributed by arXiv under the Creative Commons Attribution 4.0 International licence, http://creativecommons.org/licenses/by/4.0/, as recorded in the arXiv licence metadata for this exact version and shown on the abstract page https://arxiv.org/abs/2505.09885v2. Full licence notice: \texttt{licenses/arxiv-2505.09885-CC-BY-4.0.txt}.\par\smallskip

\noindent\textit{Editorial scope.} Counted unit: the article's prop prop (source0.tex lines 1165--1168). The article's complete original proof of it is delivered with it (source0.tex lines 1169--1202, one \texttt{proof} environment of 34 lines). No mathematics is written, repaired or re-derived: the statement and the proof are the article's own lines, in the article's own notation, and no retained line is substituted or re-flowed.  No further source-local context is needed: every cross-reference of every retained line resolves inside the two delivered spans.  Nothing was omitted from any retained span, so the package carries no omission record.  No retained line contains a citation, so the wrapper carries no bibliography and no bibliography entry.  No retained line contains a figure environment, an image-inclusion command or drawing code, so no image is delivered and the package declares no asset dependency.  Editorial formatting only, in addition: an AMS \texttt{amsart} preamble that restates the article's own declarations and macros used by the retained lines, namely the environments \texttt{prop} on separate counters, together with the standard packages those lines need.  The retained environments print in this document as Proposition 1 in order of appearance, whereas the article prints the same items with the numbering of its own full document.  The complete native source of the article as distributed in the arXiv e-print for this exact version is bundled unchanged as \texttt{sources/178143/source0.tex}.
\begin{prop}
Let $\tilde\omega=d\tilde\chi$ and $\cK_\omega$ denotes a disribution on $E$ generated by vertical vector fields $k$ such that $i_k\omega\in \cI$. Then 
    $(VE,\Tilde{Q}, \Tilde{\omega}) \xrightarrow{\Pi} (E,Q, \cK_\omega)$ is a presymplectic gPDE over background.
\end{prop}

\begin{proof} 
As explained before, since $\cK_{\omega}$ is vertical there is a natural lift $\Tilde{\cK}$ of $\cK_\omega$. Namely, $\Tilde{\cK}$ is generated by the natural lifts $\tilde k$ of  vector fields $k \in \cK_\omega$. Note that $\commut{V}{\tilde k}$ is vertical on $\tilde E$.  The statement ammounts to checking that the defining relations are satisfied, i.e. 
\begin{equation}
    \begin{gathered}
        i_{\Tilde{Q}}\Tilde{\omega} + d\Tilde{\cL} \in \cI_{\cK}\,, \\
        \half i_{\Tilde{Q}}i_{\Tilde{Q}}\Tilde{\omega} + \Tilde{Q}\Tilde{\cL} = 0 \,,
        \\
        i_{\Tilde{k}}\Tilde{\omega} \in \cI_{\cK}
\end{gathered}
\end{equation}
where $\cI_{\cK} = <\Pi^*\alpha>, \quad \alpha\in \bigwedge^1E: i_k\alpha = 0, \quad \forall k\in \cK_\omega $.

In order to check the first condition we calculate:
\begin{multline}
   \quad i_{\Tilde{Q}}\Tilde{\omega} = 
    \half i_{\Tilde{Q}}L_VL_V \Pi^* \omega = \\ \half(L_V L_V i_{\Tilde{Q}} \Pi^*\omega +(i_{[\Tilde{Q},V]})L_V \Pi^*\omega +  L_V i_{[\Tilde{Q},V]}\Pi^*\omega) \,.\quad 
\end{multline}
The first term equals $d(L_VL_V\Pi^*(H)) + \cI_{\cK}$ because $i_Q \omega = dH + \alpha, \alpha \in \cI$ and $\alpha$ is trivially annihilated by $\cK_\omega$ since $\cK_\omega$ is vertical.  
The last term is zero because $[\Tilde{Q},V]$ is $\Pi$-vertical and the second equals to $i_{[[\Tilde{Q},V],V]}\Pi^*\omega$ and hence belongs  $\cI_K$. 

That $\half i_{\Tilde{Q}}i_{\Tilde{Q}}\Tilde{\omega} + \Tilde{Q}\Tilde{\cL} = 0$ is a matter of direct check. Namely:
\begin{multline}
    i_{\Tilde{Q}}i_{\Tilde{Q}}\Tilde{\omega}  = \half i_{\Tilde{Q}}i_{\Tilde{Q}}L_V L_V \Pi^*\omega = \\= \half i_{\Tilde{Q}}( L_V L_V i_{\Tilde{Q}} \Pi^*\omega + i_{[[\Tilde{Q},V],V]}\Pi^*\omega) = \\ =  \half (L_V i_{\Tilde{Q}} L_V i_{\Tilde{Q}} \Pi^*\omega -i_{[[\Tilde{Q},V]} L_V i_{\Tilde{Q}} \Pi^*\omega + i_{[[\Tilde{Q},V],V]} i_{\Tilde{Q}}\Pi^*\omega ) = \\ = \half(L_V L_V i_{\Tilde{Q}}i_{\Tilde{Q}} \Pi^*\omega  + 2i_{[[\Tilde{Q},V],V]} i_{\Tilde{Q}}\Pi^*\omega )  = \\ = \half (L_V L_V i_{\Tilde{Q}}i_{\Tilde{Q}} \Pi^*\omega  - 2L_{[[\Tilde{Q},V],V]}\Pi^*\cL) \,,
\end{multline}
where we again extensively used the fact that  $i_{[\Tilde{Q},V]}\Pi^*\omega$ is zero because $[\Tilde{Q},V]$ is $\Pi$ vertical. We also have 
\begin{equation}
    \tilde{Q}\tilde{\cL} = \half L_{\Tilde{Q}}L_V L_V \Pi^*\cL = \half (L_V L_V L_{\Tilde{Q}}\Pi^*\cL + L_{[[\Tilde{Q},V],V]}\Pi^*\cL) \,.
\end{equation}

Finally, the last condition  $i_{\Tilde{k}}\Tilde{\omega} \in \cI_{\cK}$ holds because
\begin{equation}
    i_{\Tilde{k}}\Tilde{\omega} = i_{[[\Tilde{k},V],V]}\Pi^*\omega \in \cI_{\cK}\,.
\end{equation}
\end{proof}

\end{document}
